\scanpage{284}
\mspage{234}
\diagram{\begin{tikzcd}[ampersand replacement=\&]
 U \arrow[r,hook,"{i\in Dw}"] \arrow[d,hook,"{Dw}"'] \& X \arrow[d,hook] \\
 U'_0 \arrow[r,hook,"{Dw}"] \arrow[d,"{p'}"',"{\in W}"] \& X'_0 \arrow[d,"{q'\in W}"] \\
 U' \arrow[r,hook] \& X',
\end{tikzcd}}
où le premier carré est formé d'imm. \unsure{ouvertes|ouverte}. On sait \inserted{\unsure{alors|donc}} qu'il est \inserted{$W$-}cocartésien. \unsure{Ds|Dans} le deuxième, on a $p'\in W$, donc $q'\in W$ puisque $i$ est une imm. de Dwyer donc une $W$-cofibration. Donc il est égale\deleted{\unsure{ment}} $W$-cocartésien. On conclut alors grâce à \deleted{\unsure{b}}\unsure{c)} ci-dessus.

Prouvons b, pour
\diagram{\begin{tikzcd}[ampersand replacement=\&]
 X_0 \arrow[r,"{i_1}"] \arrow[d,"{i_2\in W}"'] \& X_1 \arrow[d,"{f_1\in W}"] \\
 X_2 \arrow[r,dash,"{f_2}"'] \& \deleted{\unsure{X}}Y
\end{tikzcd}}
, on l'insère dans
\diagram{\begin{tikzcd}[ampersand replacement=\&]
 X_0 \arrow[r,"{i_1}"] \arrow[d,"{i_2\in W}"'] \& X_1 \arrow[d,"{\unsure{j_1|i_1}}"] \arrow[ddr,bend left=20,"{f_1}"] \& \\
 X_2 \arrow[r,"{j_2}"'] \arrow[drr,bend right=24,"{f_2}"'] \& \unsure{\mathscr X|\mathcal X} \arrow[dr,"{\varphi}"] \& \\
 \& \& Y
\end{tikzcd}}
, où $\mathscr X=$
\diagram{\begin{tikzcd}[ampersand replacement=\&,row sep=tiny]
 \& \& X_1\\
 \displaystyle\int \& X_0 \arrow[ur] \arrow[dr] \&\\
 \& \& X_2
\end{tikzcd}}
. Comme $i_2\in W$, on sait qu'il en résulte $j_1\in W$, et comme $f_1\in W$, il en résulte $\varphi\in W$, \unsure{q.e.d.}

Reste à prouver c). Je vais le déduire d'une

\uline{Proposition} \uline{17}. Soit un carré commutatif
\diagram{\begin{tikzcd}[ampersand replacement=\&]
 \& X_0 \arrow[r,"{i_1}"] \arrow[d,"{i_2}"'] \& X_1 \arrow[d,"{f_1}"] \\
 Q={} \& X_2 \arrow[r,dash,"{f_2}"'] \& Y
\end{tikzcd}}
dans $Cat$, et $\mathscr X_Q$ la catégorie fibrée sur
\diagram{\begin{tikzcd}[ampersand replacement=\&]
 \& 0 \& 1 \arrow[l] \\
 Q^{\sim}(\simeq\Delta^1\times\Delta^1)={} \& 2 \arrow[u] \& 3 \arrow[l] \arrow[u]
\end{tikzcd}}
qui l'intègre. Soit $\mathscr X$ l'image inverse du fermé
\diagram{\begin{tikzcd}[ampersand replacement=\&,row sep=small,column sep=small]
 0 \& 1 \arrow[l] \\
 2 \arrow[u] \&
\end{tikzcd}}
% archive sheet number: 284

